\documentclass[11pt,a4paper]{article}

\usepackage[margin=2.2cm]{geometry}
\usepackage[T1]{fontenc}
\usepackage[utf8]{inputenc}
\usepackage[portuguese]{babel}
\usepackage{lmodern}
\usepackage{microtype}
\usepackage{array}
\usepackage{amsmath}
\usepackage{amssymb}
\usepackage{booktabs}
\usepackage{longtable}
\usepackage{xcolor}
\usepackage[hidelinks]{hyperref}

\setlength{\parindent}{0pt}
\setlength{\parskip}{0.55em}
\newcommand{\validar}{\textcolor{blue!70!black}{\textbf{A validar}}}
\newcommand{\dura}{\textbf{Restrição obrigatória}}
\newcommand{\suave}{\textbf{Preferência otimizada}}

\title{Escalas dos Técnicos de Radiologia\\
Documento de Restrições para Validação pela Equipa Hospitalar}
\author{Projeto de apoio à decisão para elaboração de escalas}
\date{Versão de \today}

\begin{document}
\maketitle

\section{Objetivo deste documento}

Este documento descreve, em linguagem acessível, as restrições e preferências
atualmente implementadas no protótipo computacional de elaboração mensal das
escalas dos técnicos de radiologia. O seu objetivo é permitir que a equipa do
Hospital da Covilhã confirme, corrija ou complete as regras antes de estas serem
consideradas representativas da prática real.

O sistema utiliza programação por restrições para procurar uma escala que
satisfaça todas as regras obrigatórias e, entre as escalas admissíveis, procure
uma distribuição equilibrada do trabalho e das tarefas menos desejáveis.

\textbf{Nota importante:} as regras abaixo descrevem o que está implementado no
software. Não constituem, por si só, uma interpretação validada da legislação,
dos contratos de trabalho, das normas clínicas ou dos procedimentos internos do
Hospital. Esses elementos devem ser confirmados pelas pessoas competentes.

\section{Terminologia}

\begin{longtable}{>{\raggedright\arraybackslash}p{0.18\textwidth}
                  >{\raggedright\arraybackslash}p{0.24\textwidth}
                  >{\raggedright\arraybackslash}p{0.48\textwidth}}
\toprule
\textbf{Código} & \textbf{Horário implementado} & \textbf{Descrição} \\
\midrule
\endhead
M  & 08:00--16:00 & Turno da manhã. \\
A  & 16:00--00:00 & Turno normal da tarde. \\
EA & 14:00--22:00 & Turno antecipado da tarde. Para efeitos de cobertura, conta como turno da tarde. \\
N  & 00:00--08:00 & Turno da noite. \\
\bottomrule
\end{longtable}

Todos os turnos são contabilizados como oito horas de trabalho.

\section{Restrições obrigatórias atualmente implementadas}

Uma restrição obrigatória não pode ser violada. Se todas as restrições
obrigatórias forem incompatíveis, o sistema não produz uma escala.

\subsection{Uma atribuição diária por técnico}

Cada técnico pode realizar, no máximo, um turno por dia civil. Um técnico pode
também ficar sem turno nesse dia.

\textbf{Questão para validação:} \validar{} --- existem situações em que um
técnico realiza dois períodos de trabalho no mesmo dia, mesmo que a soma não
ultrapasse o horário contratual?

\subsection{Cobertura mínima por turno}

\begin{center}
\begin{tabular}{lccc}
\toprule
& \textbf{Manhã} & \textbf{Tarde (A ou EA)} & \textbf{Noite} \\
\midrule
Segunda a sexta-feira & 7 & 3 & 1 \\
Sábado e domingo      & 3 & 3 & 1 \\
\bottomrule
\end{tabular}
\end{center}

Estes valores são mínimos. O modelo permite, tecnicamente, colocar mais pessoas
do que o mínimo, embora a solução atual tenha exatamente o número mínimo de
atribuições necessário.

\textbf{Questões para validação:}
\begin{itemize}
    \item \validar{} --- estes números são corretos para todos os dias do mês?
    \item \validar{} --- feriados devem seguir a regra de fim de semana ou uma regra própria?
    \item \validar{} --- existem dias, serviços ou modalidades com necessidades diferentes?
    \item \validar{} --- deve existir também um limite máximo de técnicos por turno?
\end{itemize}

\subsection{Turno antecipado da tarde}

Em cada dia útil, exatamente uma das pessoas atribuídas à tarde é marcada como
EA, trabalhando das 14:00 às 22:00. Ao sábado e ao domingo não existe EA.

\textbf{Questão para validação:} \validar{} --- esta regra também se aplica a
feriados em dias úteis e a todos os períodos do ano?

\subsection{Férias e indisponibilidade}

Um técnico não pode ser escalado numa data registada como férias ou
indisponibilidade. Atualmente, férias e indisponibilidades são representadas da
mesma forma no modelo.

\textbf{Questões para validação:}
\begin{itemize}
    \item \validar{} --- é necessário distinguir férias, baixa, formação e outras ausências?
    \item \validar{} --- existem indisponibilidades parciais que impeçam apenas alguns turnos?
\end{itemize}

\subsection{Elegibilidade individual para turnos}

As seguintes regras encontram-se diretamente codificadas:

\begin{longtable}{>{\raggedright\arraybackslash}p{0.24\textwidth}
                  >{\raggedright\arraybackslash}p{0.66\textwidth}}
\toprule
\textbf{Técnico} & \textbf{Turnos permitidos atualmente} \\
\midrule
\endhead
Celia L & Apenas manhã, em qualquer dia da semana. \\
Fernanda & Manhã ou tarde, em qualquer dia da semana; não realiza noites. \\
Sandra & Manhã ou tarde, em qualquer dia da semana; não realiza noites. \\
Celso & Manhã ou tarde, apenas de segunda a sexta-feira. \\
Angelo & Apenas manhã, de segunda a sexta-feira. \\
Nuno & Apenas manhã, de segunda a sexta-feira. \\
Telmo, Anabela, Raquel B, Rosa R, Pedro, Lina, Hugo, Richard, Vanessa, Raquel M e Novo C
& Manhã, tarde ou noite, em qualquer dia da semana. \\
\bottomrule
\end{longtable}

\textbf{Questão para validação:} \validar{} --- estas limitações refletem
qualificações, contratos, condições temporárias ou preferências? Deve ser
registada a respetiva justificação e data de validade?

\subsection{Limite semanal de horas}

Cada técnico pode trabalhar, no máximo, 40 horas por semana ISO (segunda-feira
a domingo), ou seja, no máximo cinco turnos de oito horas. Nas semanas que
começam ou terminam fora do mês, o modelo apenas conhece os dias pertencentes ao
mês que está a planear, salvo informação adicional sobre a escala anterior.

\textbf{Questões para validação:}
\begin{itemize}
    \item \validar{} --- o limite correto é 40 horas para todos os técnicos?
    \item \validar{} --- o limite deve considerar semanas móveis, médias de várias semanas ou trabalho realizado no mês anterior?
\end{itemize}

\subsection{Limite mensal para contratos de 35 ou 40 horas}

Cada técnico tem agora uma carga semanal configurável de 35 ou 40 horas. O
limite mensal é proporcional ao número de dias do mês e arredondado por defeito
para turnos completos de oito horas. Subtraem-se depois as horas de descanso
transitadas. Para um contrato de 35 horas, este cálculo introduz aproximadamente
uma folga adicional por cada oito turnos relativamente a um ciclo de 40 horas.

\textbf{Assunção provisória:} todos os técnicos foram configurados com 35 horas,
por ainda não estar disponível a carga individual. \validar{} --- indicar quem
tem contrato de 35 horas e quem tem contrato de 40 horas, e confirmar a regra de
arredondamento e transporte entre meses.

\subsection{Descanso entre tarde e manhã}

Um turno normal da tarde (16:00--00:00) não pode ser seguido por um turno da
manhã no dia seguinte. A combinação EA (14:00--22:00) seguida de manhã é
permitida.

\textbf{Questão para validação:} \validar{} --- o intervalo EA--M é aceite pelas
regras legais e internas aplicáveis?

\subsection{Máximo de tardes consecutivas}

\textbf{Nova regra confirmada pela equipa hospitalar:} nenhum técnico pode
realizar mais de três turnos de tarde consecutivos (turno normal ou EA, em
qualquer combinação). Este limite de três é obrigatório; o objetivo preferido
é não passar de duas tardes seguidas, tratado como preferência na secção
seguinte.

\subsection{Descanso após turno noturno}

Depois de um turno N, o técnico não pode realizar qualquer turno no dia civil
seguinte.

Se for fornecida a escala do mês anterior, um turno N no último dia desse mês
impede trabalho no primeiro dia do novo mês. Um turno A no último dia do mês
anterior impede apenas M no primeiro dia; um EA anterior ainda não é tratado de
forma distinta nesta fronteira.

\textbf{Questões para validação:}
\begin{itemize}
    \item \validar{} --- um dia completo sem trabalho é o descanso correto após uma noite?
    \item \validar{} --- devem ser verificadas horas exatas de descanso em vez de dias civis?
    \item \validar{} --- que regras devem ser aplicadas entre o último dia do mês planeado e o mês seguinte?
\end{itemize}

\subsection{Espaçamento entre noites}

Para o mesmo técnico, os inícios de dois turnos noturnos devem ficar separados
por pelo menos seis dias de calendário. Assim, uma noite no dia 1 impede outra
noite nos dias 2 a 6, sendo novamente possível no dia 7.

\textbf{Questão para validação:} \validar{} --- este intervalo é o utilizado na prática?

\subsection{Distribuição obrigatória das noites}

O número mínimo de noites do mês é dividido pelo número de técnicos elegíveis
para noites. Cada técnico elegível deve receber um número de noites entre o
arredondamento por defeito e o arredondamento por excesso dessa média. Para 31
noites e 11 técnicos elegíveis, cada um recebe obrigatoriamente duas ou três
noites.

Esta regra trata a distribuição das noites como uma restrição obrigatória, e não
apenas como uma preferência. Pode causar inviabilidade quando um técnico está
muito tempo indisponível.

\textbf{Questão para validação prioritária:} \validar{} --- a igualdade das
noites deve ser obrigatória ou apenas otimizada, tendo em conta férias,
contratos, qualificações e preferências?

\subsection{Regra especial de cobertura por Lina}

Quando Angelo ou Nuno estão sem trabalhar, Lina deve realizar o turno da manhã,
exceto quando Lina está de férias. Na implementação atual, esta regra também se
aplica aos fins de semana. Como Angelo e Nuno só podem trabalhar manhãs em dias
úteis, Lina é forçada a trabalhar de manhã ao sábado e ao domingo, salvo férias.

\textbf{Questão para validação prioritária:} \validar{} --- esta regra está
correta? Deve aplicar-se apenas em dias úteis, apenas em determinadas funções ou
apenas quando ambos estão ausentes?

\subsection{Prevenção de Angelo e Nuno}

Em cada dia do calendário existe exatamente uma prevenção, realizada por Angelo
ou Nuno. A prevenção é independente do turno normal, pelo que pode surgir na
escala como P ou combinada com a manhã como M/P.

\textbf{Horário confirmado pela equipa hospitalar:} a prevenção decorre das
16:00 às 08:00 do dia seguinte nos dias úteis, e das 08:00 às 08:00 do dia
seguinte (24 horas) ao sábado e ao domingo. O modelo não contabiliza estas
horas de prevenção como horas de turno; a prevenção é apenas registada por dia
civil.

\textbf{Atualização confirmada pela equipa hospitalar:} de segunda a
quinta-feira, a pessoa de prevenção alterna obrigatoriamente todo os dias. De
sexta a domingo, a prevenção passou a ser tratada como um único bloco rotativo
de três dias: a mesma pessoa assegura sexta, sábado e domingo, e essa
responsabilidade roda entre Angelo e Nuno de fim de semana para fim de semana.
Nas fronteiras do bloco (quinta para sexta e domingo para a segunda seguinte) a
pessoa de prevenção continua obrigatoriamente a mudar. Esta correção substitui a
versão anterior do modelo, que alternava a prevenção todos os dias sem
distinguir o fim de semana.

Uma prevenção realizada ao domingo ou num feriado gera uma folga compensatória
distinta, identificada por FC num dia útil sem turno regular. A implementação
valida que todas essas prevenções têm uma FC associada e que a folga não coincide
com férias nem com trabalho regular.

\textbf{Questão para validação:} \validar{} --- quando um feriado cai de
segunda a quinta-feira, deve manter-se a alternância diária ou deve juntar-se
ao bloco rotativo do fim de semana mais próximo?

\section{Preferências atualmente otimizadas}

As regras desta secção podem ser violadas quando necessário. Cada violação tem
um custo e o sistema procura minimizar a soma ponderada desses custos.

\subsection{Equilíbrio da carga de trabalho}

O objetivo principal aproxima as horas de cada técnico de uma quota proporcional
aos seus dias disponíveis. A disponibilidade é atualmente o número de dias em
que a pessoa não está de férias e pode realizar pelo menos um tipo de turno.

Este cálculo não utiliza ainda horas contratuais individuais, percentagem de
tempo, qualificações por serviço nem preferências declaradas.

\subsection{Ciclo de trabalho e agrupamento das folgas}

O modelo procura evitar mais de cinco dias consecutivos de trabalho, aproximando
o ciclo habitual de cinco turnos e duas folgas. A regra anterior, que penalizava
mais de duas folgas seguidas, foi removida. O sistema penaliza agora uma única
folga entre dois dias de trabalho e também um único turno entre duas folgas.
Assim, favorece blocos de duas ou mais folgas e evita turnos isolados.

As férias interrompem a avaliação destes padrões. \validar{} --- confirmar se o
limite de cinco dias deve ser obrigatório ou preferencial e como tratar padrões
que atravessam a fronteira entre meses.

\subsection{Preferência por, no máximo, duas tardes seguidas}

\textbf{Nova regra confirmada pela equipa hospitalar:} o sistema penaliza um
terceiro turno de tarde consecutivo (turno normal ou EA). Duas tardes seguidas
não têm penalização; a partir da terceira, aplica-se o custo desta preferência.
O limite absoluto de três tardes seguidas mantém-se obrigatório, conforme
descrito na secção de restrições obrigatórias.

\subsection{Sequência de turnos antes de uma noite}

\textbf{Nova regra confirmada pela equipa hospitalar:} o sistema prefere que,
antes de um turno da noite, os três dias anteriores sigam uma de duas
sequências: tarde normal, tarde antecipada (EA) e manhã (T, T14, M); ou tarde
normal, tarde normal e folga (T, T, D). Quando nenhuma das duas sequências é
possível ou escolhida, aplica-se uma penalização. Esta regra é apenas uma
preferência --- pode ser violada quando necessário para garantir cobertura ou
por outras restrições.

\textbf{Questão para validação:} \validar{} --- estas são as duas únicas
sequências aceitáveis antes de uma noite, incluindo casos junto à fronteira do
mês (quando os três dias anteriores não estão todos no mês planeado)?

\subsection{Prevenção preferencialmente fora dos dias de folga}

\textbf{Nova regra confirmada pela equipa hospitalar:} sempre que possível, o
sistema evita atribuir prevenção a Angelo ou Nuno num dia em que, de outro
modo, ficariam de folga (sem turno regular). Como a alternância diária e o
bloco rotativo de fim de semana já determinam, quase sempre, quem está de
prevenção em cada dia, esta preferência atua principalmente sobre a decisão de
colocar essa pessoa a trabalhar de manhã (M) nesse dia, em vez de lhe dar
folga, quando a cobertura o permite. A prevenção em dias de férias mantém-se
proibida, como restrição obrigatória.

\subsection{Igualdade de dias trabalhados entre meses}

O sistema minimiza, com prioridade elevada, a diferença entre os dias acumulados
de trabalho de todos os técnicos. Se a igualdade não for possível num mês, o
resultado é guardado e usado no mês seguinte, favorecendo quem ficou com menos
dias. Para maio de 2025, sem histórico anterior, a diferença mínima é um dia
porque são necessárias 305 atribuições para 17 técnicos.

\textbf{Atualização confirmada pela equipa hospitalar:} o objetivo é terminar
cada ano civil com o mesmo número de horas para todos os técnicos. Por isso, o
transporte da diferença acumulada de dias trabalhados entre meses é agora
reiniciado a zero em janeiro de cada ano, em vez de se acumular indefinidamente
ao longo de vários anos.

\subsection{Rotação de feriados e datas especiais}

O trabalho em feriados é equilibrado usando o histórico acumulado. Natal, Ano
Novo e domingo de Páscoa possuem contagens históricas separadas, para que a
rotação de cada data especial seja feita entre anos. Em maio de 2025 foi
registado o dia 1 de maio; os históricos anteriores foram inicializados a zero.

\subsection{Equilíbrio dos fins de semana}

Um fim de semana é contado como trabalhado se o técnico trabalhar no sábado,
no domingo ou em ambos. O modelo procura aproximar a razão

\[
\frac{\text{fins de semana trabalhados}}
     {\text{fins de semana em que poderia trabalhar}}
\]

entre os técnicos elegíveis.

\subsection{Fins de semana consecutivos}

Trabalhar em dois fins de semana consecutivos gera uma penalização, mas continua
a ser permitido quando necessário para garantir cobertura.

\subsection{Fim de semana incompleto}

Trabalhar apenas no sábado ou apenas no domingo gera uma penalização. O modelo
prefere, quando possível, que uma pessoa trabalhe os dois dias ou nenhum deles.

\subsection{Fins de semana próximos de férias}

O modelo penaliza trabalho em fins de semana associados a períodos de férias.
Para cada dia de férias, são considerados o fim de semana imediatamente anterior,
o fim de semana associado à mesma semana e o fim de semana imediatamente
seguinte. A união destes períodos pode abranger vários fins de semana quando as
férias são longas.

\textbf{Questão para validação:} \validar{} --- qual é exatamente o período de
proteção desejado antes e depois das férias?

\subsection{Estabilidade após alterações}

Quando uma escala já publicada precisa de ser reparada devido a uma ausência ou
aumento da procura, o modelo penaliza cada atribuição removida ou acrescentada.
Nesta modalidade, preservar a escala original recebe um peso elevado. São
distinguidos M, A, EA e N.

\section{Pesos utilizados na escala normal}

Na configuração normal atual, os componentes da função objetivo têm os pesos
seguintes:

\begin{center}
\begin{tabular}{lr}
\toprule
\textbf{Componente} & \textbf{Peso} \\
\midrule
Diferença acumulada de dias trabalhados & 100,000 \\
Desvio máximo da carga proporcional & 100 \\
Mais de cinco dias consecutivos de trabalho & 500 \\
Turno isolado entre folgas & 300 \\
Folga isolada entre turnos & 200 \\
Terceira tarde consecutiva & 400 \\
Sequência de turnos antes de uma noite não seguida & 150 \\
Prevenção coincidente com dia de folga & 250 \\
Rotação de feriados & 1000 \\
Rotação de Natal, Ano Novo e Páscoa & 2000 \\
Desequilíbrio normalizado de fins de semana & 30 \\
Fins de semana consecutivos & 20 \\
Fim de semana incompleto & 10 \\
Fim de semana próximo de férias & 5 \\
Alterações à escala publicada & 0 na criação inicial; 100,000 na reparação \\
\bottomrule
\end{tabular}
\end{center}

Estes pesos implementam uma soma ponderada, não uma hierarquia matemática
estritamente lexicográfica. Uma melhoria grande num componente pode compensar
várias pequenas perdas noutro.

\textbf{Questão para validação:} \validar{} --- qual deve ser a ordem de
prioridade entre equilíbrio de horas, descanso, noites, fins de semana,
preferências pessoais e estabilidade da escala?

\section{Aspetos ainda não representados}

Os seguintes elementos referidos ou sugeridos no projeto ainda não estão
representados de forma completa no modelo:

\begin{itemize}
    \item confirmação das cargas contratuais individuais de 35 ou 40 horas;
    \item qualificações por modalidade, equipamento, sala ou função clínica;
    \item níveis mínimos de experiência por turno;
    \item preferências declaradas individualmente pelos técnicos;
    \item pedidos de folga com diferentes níveis de prioridade;
    \item calendário completo dos feriados locais e nacionais;
    \item baixas, formação e outros tipos de ausência distintos de férias;
    \item horas extraordinárias, custos ou compensações;
    \item composição das equipas por grupo operacional;
    \item regras completas na fronteira entre meses;
    \item alterações temporárias de elegibilidade;
    \item validação jurídica e institucional formal de cada regra.
\end{itemize}

\section{Tabela de validação pela equipa hospitalar}

Solicita-se que cada regra seja classificada como correta, incorreta ou
incompleta. As correções podem ser registadas na última coluna.

\renewcommand{\arraystretch}{1.35}
\small
\begin{longtable}{>{\raggedright\arraybackslash}p{0.27\textwidth}
                  >{\centering\arraybackslash}p{0.10\textwidth}
                  >{\centering\arraybackslash}p{0.11\textwidth}
                  >{\centering\arraybackslash}p{0.12\textwidth}
                  >{\raggedright\arraybackslash}p{0.20\textwidth}}
\toprule
\textbf{Regra} & \textbf{Correta} & \textbf{Incorreta} & \textbf{Incompleta} & \textbf{\shortstack{Correção/\\observação}} \\
\midrule
\endhead
Um turno por técnico e por dia & $\square$ & $\square$ & $\square$ & \\
Cobertura em dias úteis: 7 M, 3 A, 1 N & $\square$ & $\square$ & $\square$ & \\
Cobertura ao fim de semana: 3 M, 3 A, 1 N & $\square$ & $\square$ & $\square$ & \\
Um EA por dia útil & $\square$ & $\square$ & $\square$ & \\
Elegibilidade individual & $\square$ & $\square$ & $\square$ & \\
Máximo de 40 horas por semana & $\square$ & $\square$ & $\square$ & \\
Fórmula do limite mensal & $\square$ & $\square$ & $\square$ & \\
Descanso A--M e EA--M & $\square$ & $\square$ & $\square$ & \\
Dia sem trabalho após N & $\square$ & $\square$ & $\square$ & \\
Seis dias entre inícios de noites & $\square$ & $\square$ & $\square$ & \\
Duas ou três noites por técnico elegível & $\square$ & $\square$ & $\square$ & \\
Máximo de três tardes seguidas (preferência até duas) & $\square$ & $\square$ & $\square$ & \\
Sequência T/T14/M ou T/T/D antes de uma noite & $\square$ & $\square$ & $\square$ & \\
Horário da prevenção (16--08 / 08--08 fim de semana) & $\square$ & $\square$ & $\square$ & \\
Prevenção preferencialmente fora dos dias de folga & $\square$ & $\square$ & $\square$ & \\
Regra de Lina, Angelo e Nuno & $\square$ & $\square$ & $\square$ & \\
Prevenção: bloco rotativo sexta--domingo & $\square$ & $\square$ & $\square$ & \\
Igualdade de horas com reinício anual em janeiro & $\square$ & $\square$ & $\square$ & \\
Ciclo de cinco turnos e duas folgas & $\square$ & $\square$ & $\square$ & \\
Folgas agrupadas e ausência de turno isolado & $\square$ & $\square$ & $\square$ & \\
Compensação de dias no mês seguinte & $\square$ & $\square$ & $\square$ & \\
Rotação de feriados e datas especiais & $\square$ & $\square$ & $\square$ & \\
Equilíbrio e continuidade dos fins de semana & $\square$ & $\square$ & $\square$ & \\
Proteção de fins de semana junto às férias & $\square$ & $\square$ & $\square$ & \\
Prioridade dos objetivos & $\square$ & $\square$ & $\square$ & \\
\bottomrule
\end{longtable}
\normalsize

\section{Confirmação}

Nome e função: \hrulefill

Data: \hrulefill

Observações gerais:

\vspace{3cm}
\hrulefill

\vspace{1.2cm}
\hrulefill

\end{document}
